% =============================================================
% Domain-Informed Loss for Structured Sequence Prediction:
% Articulatory Distance in Grapheme-to-Phoneme Conversion
% IEEE MLSP 2026 — Double-blind submission (6 pages incl. refs)
% =============================================================
%
% COMPILATION: latexmk -pdf main
% =============================================================

\documentclass[conference]{IEEEtran}
\IEEEoverridecommandlockouts

% ---- Encoding & fonts ----------------------------------------
\usepackage[utf8]{inputenc}
\usepackage[T1]{fontenc}
\usepackage{tipa}

% ---- Math ----------------------------------------------------
\usepackage{amsmath}
\usepackage{amssymb}

% ---- Tables --------------------------------------------------
\usepackage{booktabs}
\usepackage{multirow}

% ---- References ----------------------------------------------
\usepackage{cite}

% ---- Macros --------------------------------------------------
\newcommand{\dpp}{d_{\mathrm{PP}}}        % PanPhon distance
\newcommand{\LCE}{L_{\mathrm{CE}}}        % CE loss
\newcommand{\lam}{\lambda}                % lambda shorthand

% =============================================================
\begin{document}
% =============================================================

\title{Domain-Informed Loss Functions for Structured\\
       Sequence Prediction: Articulatory Distance\\
       in Grapheme-to-Phoneme Conversion}

\author{%
  \IEEEauthorblockN{[ANONYMOUS]}
  \IEEEauthorblockA{[ANONYMOUS]}
}

\maketitle

% ---- Abstract -----------------------------------------------
\begin{abstract}
Standard cross-entropy (CE) training for sequence prediction treats
all output errors equally, discarding domain structure that
distinguishes near-miss errors from catastrophic ones.
We propose a general framework for \textbf{domain-informed loss
functions} that augment CE with a distance penalty scaled by
prediction confidence, and instantiate it as Distance-Aware (DA)
Loss for grapheme-to-phoneme (G2P) conversion using PanPhon
articulatory distance.

The composite training signal
$L = \LCE + \lam \cdot d(\hat{y}, y) \cdot p(\hat{y})$
encodes domain-specific output similarity via $d$ and model
confidence via $p(\hat{y})$.
Applied to Brazilian Portuguese G2P with a BiLSTM encoder-decoder,
DA~Loss reduces catastrophic (Class~D) substitutions by 19\,\%
relative while maintaining overall error rate, demonstrating that
domain-informed losses reshape the \emph{distribution} of errors
rather than merely reducing their count.

We validate on a stratified 28{,}782-word test set with Wilson
95\,\% CIs ($\pm$0.03\,pp), achieving PER~0.48\,\% and
WER~5.33\,\%.
A factorial analysis reveals non-additive interaction between loss
function, model capacity, and output representation---the best
configuration requires all three factors jointly.
The framework generalizes to any structured prediction task where
a domain-grounded output distance metric is available.
\end{abstract}

\begin{IEEEkeywords}
domain-informed loss, structured prediction, distance-aware
training, grapheme-to-phoneme, articulatory features
\end{IEEEkeywords}

% =============================================================
\section{Introduction}
\label{sec:intro}
% =============================================================

Cross-entropy (CE) is the default training objective for
classification and sequence prediction.
For structured output spaces---where some errors are ``closer''
to the target than others---CE provides no gradient signal about
error severity: a near-miss and a catastrophic error receive
identical penalties when $p_\mathrm{correct}$ is equal.

This limitation is well-known in computer vision, where label
smoothing~\cite{szegedy2016rethinking} redistributes probability
mass uniformly---ignoring output structure---and in speech
recognition, where edit-distance losses operate at the sequence
level.
The specific formulation we propose---combining a domain distance
metric with prediction confidence as a multiplicative signal
alongside CE---has not been systematically studied for token-level
structured prediction.

We instantiate this framework in grapheme-to-phoneme (G2P)
conversion, where the output space has a natural metric: PanPhon
articulatory distance~\cite{mortensen2016panphon}, encoding 24
binary features (voicing, nasality, place, manner) that define how
``far apart'' two phonemes are in articulatory space.
Predicting /\textipa{E}/ when the target is /e/ ($d\!=\!0.10$, one
feature difference) is qualitatively different from predicting /k/
for /a/ ($d\!=\!0.90$, eight features)---yet CE treats them
identically.

\textbf{Comparison with LatPhon~\cite{chary2025latphon}.}
LatPhon is a 4-layer multilingual Transformer (7.5M parameters)
reporting PER~0.86\,\% (Wilson CI [0.56\,\%, 1.16\,\%]) on
$\approx$500~PT-BR words from the same dictionary used here.
Our system achieves PER~0.48\,\% (CI~[0.46\,\%, 0.51\,\%]) on
28{,}782~words---the CIs do not overlap, a statistically significant
difference at 95\,\% confidence.
The 57$\times$ test set difference confers 10$\times$ more precise
CIs ($\pm$0.03\,pp vs.\ $\pm$0.30\,pp).
We deliberately use a BiLSTM to isolate the loss contribution from
architectural novelty.

\textbf{Contributions:}
\begin{enumerate}
  \item A general \textbf{domain-informed loss} formulation
        combining output-space distance with prediction
        confidence (Sec.~\ref{sec:method})
  \item \textbf{Error quality taxonomy}: Class~A--D classification
        showing systematic redistribution of catastrophic errors
        (Sec.~\ref{sec:analysis})
  \item \textbf{Factorial analysis} of loss $\times$ capacity
        $\times$ representation interaction, demonstrating
        non-additive effects (Sec.~\ref{sec:factorial})
  \item Large-scale evaluation with \textbf{stratified splits} and
        Wilson CIs on 28{,}782~words (Sec.~\ref{sec:eval})
\end{enumerate}

% =============================================================
\section{Related Work}
\label{sec:related}
% =============================================================

\textbf{Loss function design for structured outputs.}
Label smoothing~\cite{szegedy2016rethinking} uniformly
redistributes probability mass, ignoring output structure.
Knowledge distillation uses soft targets from a teacher model but
does not incorporate domain-specific distance.
Edit-distance losses have been explored in speech recognition but
typically at the sequence level, not per-token.

\textbf{Domain knowledge in training signals.}
Phonological distance has been used in ASR for confusion-weighted
evaluation, but not as a direct training signal.
Our work differs in three aspects: (1)~we operate at the IPA token
level; (2)~we use PanPhon articulatory
distance~\cite{mortensen2016panphon} as a principled,
language-independent metric; (3)~we weight the distance penalty by
prediction confidence, creating a composite signal that maximally
penalizes errors that are both confident and distant.

\textbf{G2P systems.}
Neural G2P has evolved from RNN
approaches~\cite{rao2015g2p,bisani2008joint} to Transformer
architectures.
LatPhon~\cite{chary2025latphon} applies a multilingual Transformer
to 6~languages including Portuguese.
We deliberately use a BiLSTM~\cite{bahdanau2014neural} to isolate
the loss function contribution from architectural novelty.

% =============================================================
\section{Method: Domain-Informed Loss}
\label{sec:method}
% =============================================================

\subsection{General Framework}

For any structured prediction task with output vocabulary
$\mathcal{V}$ and a domain-grounded distance $d: \mathcal{V}
\times \mathcal{V} \to [0,1]$, we define:
\begin{equation}
  \boxed{L = \LCE + \lam \cdot d(\hat{y},\, y) \cdot p(\hat{y})}
  \label{eq:diloss}
\end{equation}
where $\LCE = -\log p_\mathrm{correct}$,
$d(\hat{y}, y)$ is the domain distance between predicted and target
outputs, $p(\hat{y})$ is the softmax probability assigned to the
(incorrect) argmax prediction, and $\lam$ controls coupling strength.

\textbf{Key properties:}
(1)~CE and domain terms monitor \emph{independent} quantities:
$p_\mathrm{correct}$ vs.\ $p(\hat{y})$.
When the model is confidently wrong, $p_\mathrm{correct}$ is low
(CE high) while $p(\hat{y})$ is high (domain penalty high)---the
two signals reinforce.
(2)~The domain term is \emph{bounded}
($\leq \lam \times 1.0 \times 1.0$), ensuring CE remains the
primary learning axis.
(3)~The framework is \textbf{task-agnostic}: any domain providing
a normalized output distance can instantiate~(\ref{eq:diloss}).

\subsection{Instantiation: Articulatory Distance for G2P}

For G2P, $d = \dpp$: normalized Euclidean distance in PanPhon's
24-D articulatory feature space~\cite{mortensen2016panphon}.
Each phoneme is encoded as a 24-dimensional vector of binary
articulatory features (voicing, nasality, place, manner).
Example distances:
p$\leftrightarrow$b: 0.04 (voicing only);
s$\leftrightarrow$\textipa{S}: 0.15 (place shift);
a$\leftrightarrow$k: 0.90 (vowel vs.\ velar stop).

\textbf{Structural token override.}
Non-phonemic tokens (syllable boundary ``.'', stress marker
``\textipa{"}'') receive zero vectors in PanPhon.
Without correction, $d(., \text{'}) = 0$---DA provides no gradient
signal for their confusion.
We apply a post-normalization override: $d = 1.0$ for any pair
involving a structural token.
The override must be applied \emph{after} Euclidean normalization;
applying it before yields $d\!\approx\!0.25$, not the intended
maximum.

\textbf{Numerical example.}
Consider two errors with identical CE ($p_\mathrm{correct}\!=\!0.40$,
CE$\,=\,0.916$):
a near-miss e$\to$\textipa{E} ($d\!=\!0.10$, $p(\hat{y})\!=\!0.45$)
adds DA$\,=\,0.009$, total $L\!=\!0.925$;
a catastrophic a$\to$k ($d\!=\!0.90$, $p(\hat{y})\!=\!0.45$) adds
DA$\,=\,0.081$, total $L\!=\!0.997$.
CE alone is identical (0.916) in both cases.
DA provides a 9$\times$ difference in phonological penalty, steering
the model to ``break ties toward the closer candidate.''

\textbf{Worked example.}
Consider the word \textit{cena} (target: /s\,e\,n\,\textipa{@}/)
at the decoder step for grapheme ``c''.
The softmax produces
$\{$s:0.42, \textipa{S}:0.35, z:0.12, t:0.06, \ldots$\}$;
$\hat{y}\!=\!\textipa{S}$ (argmax, incorrect).
CE: $-\!\log(0.42)\!=\!0.868$ (penalizes low $p_\mathrm{correct}$).
DA: $0.20\!\times\!d(\textipa{S},\mathrm{s})\!\times\!p(\textipa{S})
     = 0.20\!\times\!0.15\!\times\!0.35 = 0.0105$.
Total: $L\!=\!0.879$.
If the model had predicted /k/ instead:
DA$\,=\,0.20\!\times\!0.80\!\times\!0.35 = 0.056$;
total $L\!=\!0.924$---5\,\% higher for the distant error.
Over thousands of gradient steps, this differential signal steers
the model to ``break ties toward alveolar fricatives'' without
receiving an explicit phonological rule.

\subsection{$\lambda$ Selection}

Table~\ref{tab:lambda} sweeps $\lam$ with fixed architecture
(4.3M params).
The inverted-U pattern is consistent with bounded auxiliary losses:
too low is inert; too high competes with CE.
\begin{table}[ht]
  \caption{$\lambda$ sweep, 4.3M architecture.}
  \label{tab:lambda}
  \centering\small
  \begin{tabular}{cccc}
    \toprule
    $\lambda$ & PER & WER & Behavior \\
    \midrule
    0.05 & 0.62\,\% & 5.36\,\% & Signal too weak \\
    0.10 & 0.63\,\% & 5.35\,\% & Moderate \\
    \textbf{0.20} & \textbf{0.60\,\%} & \textbf{5.14\,\%} & \textbf{Optimal} \\
    0.50 & 0.65\,\% & 5.57\,\% & Over-penalization \\
    \bottomrule
  \end{tabular}
\end{table}

% =============================================================
\section{Experimental Setup}
\label{sec:eval}
% =============================================================

\subsection{Task and Data}

We evaluate on Brazilian Portuguese G2P: mapping orthographic words
to IPA phoneme sequences.
The corpus contains \textbf{95{,}937 word--IPA pairs}, split
\textbf{60/10/30} (train/val/test) with stratification over stress
type (oxytone/paroxytone/proparoxytone), syllable count (1--5+),
and character length ($\leq$4, 5--7, 8--10, 11+), yielding
$\sim$48 strata.
Split quality: $\chi^2\!=\!0.95$ ($p\!=\!0.678$),
Cram\'{e}r $V\!=\!0.0007$~\cite{kohavi1995crossvalidation}.

\textbf{Split bias evidence.}
An unstratified 70/10/20 split (same architecture, CE loss) achieved
PER~1.12\,\%, while the stratified 60/10/30 split achieved
0.66\,\%---a 41\,\% PER reduction attributable entirely to
evaluation protocol, not model improvement.

\subsection{Architecture}

BiLSTM encoder-decoder with Bahdanau
attention~\cite{bahdanau2014neural}, chosen to isolate the loss
contribution from architectural novelty.

\textbf{Encoder}: 2-layer BiLSTM over grapheme embeddings.
Each position $t$ yields
$h_t\!=\![\overrightarrow{h}_t;\,\overleftarrow{h}_t]$,
providing bidirectional context---critical for resolving graphemic
ambiguity (e.g., ``c'' $\to$ /k/ or /s/ depends on following vowel).

\textbf{Decoder}: 2-layer LSTM with Bahdanau attention and teacher
forcing during training; autoregressive at inference.

Three capacity configurations (Table~\ref{tab:configs}).

\begin{table}[ht]
  \caption{Architecture configurations.}
  \label{tab:configs}
  \centering\small
  \begin{tabular}{lccc}
    \toprule
    Config  & Embedding & Hidden & Parameters \\
    \midrule
    Small   & 128-D & 256-D & 4.3M \\
    Medium  & 192-D & 384-D & 9.7M \\
    Large   & 256-D & 512-D & 17.2M \\
    \bottomrule
  \end{tabular}
\end{table}

All models: Adam ($\beta_1\!=\!0.9$, $\beta_2\!=\!0.999$),
lr $10^{-3}$, batch size 96, early stopping
(patience 10), seed 42, teacher forcing ratio 1.0.
Batch size 96 ensures $\geq$3 samples per phonological stratum
per batch (with $\sim$34 strata), reducing loss curve variance
by $\sim$50\,\% relative to batch size 32.
Models converge in 80--120 epochs ($\sim$59K gradient steps).

\textbf{Two levels of randomness.}
A frequently conflated distinction: \emph{stratified splitting}
(classical sampling theory) guarantees proportional representation
of phonological strata in each subset, preventing metric inflation
from unrepresentative test sets.
\emph{Training shuffling} (stochastic optimization) reshuffles
examples each epoch for gradient variance reduction.
The two are independent: stratification ensures unbiased
\emph{measurement}; shuffling ensures efficient
\emph{optimization}~\cite{bottou2010large}.

\subsection{Metrics}
\label{sec:metrics}

\textbf{PER}: phoneme-level edit distance normalized by reference
length~\cite{bisani2008joint}.
\textbf{WER}: word-level exact match.
\textbf{Wilson 95\,\% CI}~\cite{wilson1927probable}: $\pm$0.03\,pp
for our test set ($\sim$181K phonemes) vs.\ $\pm$0.30\,pp for
$\sim$500-word evaluations.
\textbf{PER$_w$}: weighted PER where each substitution is weighted
by normalized Hamming distance over PanPhon's 24 features,
capturing error \emph{quality}.

\textbf{Error taxonomy.}
We define error classes by Hamming distance $d_H$ over PanPhon's
24 articulatory features (Table~\ref{tab:errorclasses}).
Evaluation uses Hamming distance, while DA~Loss training uses
Euclidean---both over the same 24 features.
Euclidean penalizes bipolar contrasts
($+1\!\leftrightarrow\!-1$, e.g., voicing) more than unspecified
features ($0\!\leftrightarrow\!\pm 1$), which is linguistically
more appropriate as a training signal.

\begin{table}[ht]
  \caption{Error class taxonomy based on normalized Hamming distance.}
  \label{tab:errorclasses}
  \centering\small
  \begin{tabular}{lcll}
    \toprule
    Class & $d_H$ & Features diff. & Example \\
    \midrule
    A & 0.000 & 0 (exact) & --- \\
    B & $\leq$0.050 & 1 (minimal pair) & p$\leftrightarrow$b \\
    C & $\leq$0.150 & 2--3 (same family) & s$\leftrightarrow$\textipa{S} \\
    D & $>$0.150 & 4+ (different) & vowel$\leftrightarrow$stop \\
    \bottomrule
  \end{tabular}
\end{table}

% =============================================================
\section{Results}
\label{sec:results}
% =============================================================

\subsection{Main Results}

Table~\ref{tab:results} shows the complete experimental progression
across 7~methodological phases.

\begin{table*}[ht]
  \caption{Complete experimental progression. Sep.\ = syllable
           boundary tokens. Dist.\ = structural distance override.
           Bold = recommended configurations.}
  \label{tab:results}
  \centering\small
  \begin{tabular}{llcccccl}
    \toprule
    Phase & Exp & Params & Loss & Sep. & PER & WER & Insight \\
    \midrule
    \multirow{2}{*}{1.~Split}
      & Exp0 & 4.3M & CE & --- & 1.12\,\% & 9.37\,\% & Baseline (70/10/20) \\
      & Exp1 & 4.3M & CE & --- & 0.66\,\% & 5.65\,\% & Stratified 60/10/30: $-41$\,\% PER \\
    \midrule
    \multirow{2}{*}{2.~Capacity}
      & Exp5 & 9.7M & CE & --- & 0.63\,\% & 5.38\,\% & Sweet spot \\
      & Exp2 & 17.2M & CE & --- & 0.60\,\% & 4.98\,\% & Diminishing returns \\
    \midrule
    \multirow{3}{*}{3.~DA Loss}
      & Exp7 & 4.3M & DA $\lam\!=\!0.2$ & --- & 0.60\,\% & 5.14\,\% & $\lam$ optimized \\
      & \textbf{Exp9} & \textbf{9.7M} & \textbf{DA} $\boldsymbol{\lam\!=\!0.2}$ & \textbf{---} & \textbf{0.58\,\%} & \textbf{4.96\,\%} & \textbf{Best WER} \\
      & Exp10 & 17.2M & DA $\lam\!=\!0.2$ & --- & 0.61\,\% & 5.25\,\% & No gain without sep \\
    \midrule
    \multirow{2}{*}{4.~Separators}
      & Exp101 & 4.3M & CE & \checkmark & 0.53\,\% & 5.99\,\% & PER$\downarrow$ WER$\uparrow$ \\
      & Exp102 & 9.7M & CE & \checkmark & 0.52\,\% & 5.79\,\% & Capacity attenuates WER cost \\
    \midrule
    \multirow{2}{*}{5.~DA+Sep}
      & Exp103 & 9.7M & DA & \checkmark & 0.53\,\% & 5.73\,\% & Non-additive effects \\
      & \textbf{104d} & \textbf{17.2M} & \textbf{DA+dist} & \checkmark & \textbf{0.48\,\%} & \textbf{5.33\,\%} & \textbf{Best PER (reference)} \\
    \midrule
    \multirow{2}{*}{6.~Robustness}
      & Exp105 & 9.7M & DA+dist & \checkmark & 0.54\,\% & 5.87\,\% & 50\% data: $-10$\,\% PER \\
      & Exp106 & 9.7M & DA+dist & \checkmark & 0.58\,\% & 6.12\,\% & No hyphen: minimal impact \\
    \midrule
    \multicolumn{2}{l}{LatPhon~\cite{chary2025latphon}}
      & 7.5M & NLL & --- & 0.86\,\% & n/a & Multilingual Transformer \\
    \bottomrule
  \end{tabular}
\end{table*}

DA~Loss at 9.7M (Exp9) achieves the best WER (4.96\,\%).
With separators and distance override at 17.2M (Exp104d), the best
PER (0.48\,\%).
Wilson CIs for Exp104d [0.46\,\%, 0.51\,\%] do not overlap with
LatPhon's~\cite{chary2025latphon} [0.56\,\%, 1.16\,\%]---a
statistically significant difference at 95\,\% confidence.

\textbf{Separator trade-off.}
Syllable boundary tokens improve PER ($-17$\,\% to $-20$\,\%)
at the cost of WER ($+6$\,\% to $+8$\,\%).
Each misplaced separator counts as a full word error under
exact-match WER---a structural trade-off independent of loss.

\subsection{Factorial Analysis: Loss $\times$ Representation}
\label{sec:factorial}

Table~\ref{tab:factorial} shows a 2$\times$2 factorial at 9.7M
parameters.
\begin{table}[ht]
  \caption{2$\times$2 factorial: separators $\times$ DA Loss.
           Values are PER\,\% / WER\,\%.}
  \label{tab:factorial}
  \centering\small
  \begin{tabular}{lcc}
    \toprule
    & CE & DA $\lam\!=\!0.2$ \\
    \midrule
    No separators   & 0.63 / 5.38 & \textbf{0.58 / 4.96} \\
    With separators & 0.52 / 5.79 & 0.53 / 5.73 \\
    \bottomrule
  \end{tabular}
\end{table}

\textbf{Key finding: non-additive interaction.}
Without separators, DA improves both PER ($-8$\,\%) and WER
($-8$\,\%).
With separators, DA does not improve over CE at 9.7M---the
factors interact non-additively.
Separators provide the largest single-factor PER improvement
($-17$\,\%), but the best overall result (PER 0.48\,\%) requires
\emph{all three}: 17.2M capacity + separators + DA + distance
override.

\textbf{Capacity interaction.}
At 17.2M without separators, DA~Loss does not improve PER relative
to CE (0.61\,\% vs.\ 0.60\,\%).
At 17.2M \emph{with} separators and corrected distances, DA achieves
the best PER (0.48\,\%).
This refutes the hypothesis that domain losses interfere with
large-model memorization---the relevant variable is the interaction
between capacity, structural representation, and distance correction.

\subsection{PanPhon Embeddings $\times$ DA Loss}

A second 2$\times$2 factorial at 4.3M parameters isolates the
interaction between PanPhon-initialized embeddings and DA~Loss
(Table~\ref{tab:panphon_factorial}).

\begin{table}[ht]
  \caption{2$\times$2 factorial: PanPhon embeddings $\times$ DA Loss,
           4.3M params, no sep. Values are PER\,\% / PER$_w$\,\%.}
  \label{tab:panphon_factorial}
  \centering\small
  \begin{tabular}{lcc}
    \toprule
    & CE & DA $\lam\!=\!0.2$ \\
    \midrule
    Learned (random)     & 0.66 / 0.34 & 0.60 / 0.27 \\
    PanPhon$_T$ (24-D)   & 0.66 / 0.33 & 0.65 / --- \\
    \bottomrule
  \end{tabular}
\end{table}

Key findings:
(1)~PanPhon initialization with CE does not improve raw PER
(0.66\,\% in both cases)---the geometric inductive bias is not
reinforced by a phonologically-blind loss.
(2)~DA~Loss with learned embeddings achieves the largest improvement
($-9$\,\% PER, $-21$\,\% PER$_w$), confirming that DA operates via
an external distance lookup, not embedding geometry.
(3)~Combining PanPhon embeddings \emph{with} DA~Loss yields
\emph{worse} PER than DA alone (0.65\,\% vs.\ 0.60\,\%)---the two
phonological inductive biases compete rather than synergize.
\textbf{Design implication}: PanPhon features are most effective as
an \emph{external signal} in the loss function, not as an embedding
prior.

\subsection{Design Space Overview}

Table~\ref{tab:designspace} summarizes the explored encoding
$\times$ loss design space.
The two recommended configurations occupy different Pareto frontiers:
9.7M DA without separators (best WER, compact deployment) and
17.2M DA+dist with separators (best PER, for TTS and forced
alignment).

\begin{table}[ht]
  \caption{Design space: PER\,\% / WER\,\% by encoding and loss.}
  \label{tab:designspace}
  \centering\small
  \begin{tabular}{lccc}
    \toprule
    Encoding & CE & DA $\lam\!=\!0.2$ & DA+dist \\
    \midrule
    Raw 4.3M        & 0.66/5.65 & 0.60/5.14 & --- \\
    Raw 9.7M        & 0.63/5.38 & \textbf{0.58/4.96} & --- \\
    Raw 9.7M + sep  & 0.52/5.79 & 0.53/5.73 & 0.49/5.43 \\
    Raw 17.2M + sep & ---       & ---       & \textbf{0.48/5.33} \\
    \bottomrule
  \end{tabular}
\end{table}

% =============================================================
\section{Analysis: Error Quality Redistribution}
\label{sec:analysis}
% =============================================================

\subsection{Class Distribution}

Table~\ref{tab:classdist} quantifies the redistribution.
DA~Loss reduces Class~D errors \textbf{19\,\% relative}
(CE baseline $\to$ DA 9.7M), independently of raw PER improvement.
\begin{table}[ht]
  \caption{Error class distribution. DA reduces catastrophic (D)
           errors while maintaining near-miss (B) rate.}
  \label{tab:classdist}
  \centering\small
  \begin{tabular}{lcccc}
    \toprule
    System & PER & Cls~B & Cls~D & D/errors \\
    \midrule
    CE base.\ (4.3M)         & 0.66\,\% & 0.39\,\% & 0.54\,\% & 50.9\,\% \\
    DA (9.7M)                & 0.58\,\% & 0.36\,\% & 0.44\,\% & 48.4\,\% \\
    DA+dist (17.2M+sep.)     & 0.48\,\% & 0.29\,\% & 0.53\,\% & 47.4\,\% \\
    \bottomrule
  \end{tabular}
\end{table}

\subsection{Graduated Metrics}

Table~\ref{tab:graduated} tracks PER$_w$ (distance-weighted PER)
across experiments.
DA~Loss reduces PER$_w$ by $-21$\,\% (0.34\,$\to$\,0.27\,\%) vs.\
the CE baseline, while capacity alone (17.2M) reduces it only
$-15$\,\%.
This confirms that the domain-informed loss improves error
\emph{quality} disproportionately to error \emph{quantity}---the
central claim of domain-informed loss design.
\begin{table}[ht]
  \caption{Graduated metrics. PER$_w$ captures error quality.}
  \label{tab:graduated}
  \centering\small
  \begin{tabular}{llccc}
    \toprule
    System & Method & PER & PER$_w$ & WER$_g$ \\
    \midrule
    Exp1 & CE baseline       & 0.66\,\% & 0.34\,\% & 0.69\,\% \\
    Exp2 & CE, +capacity     & 0.60\,\% & 0.29\,\% & 0.62\,\% \\
    Exp6 & DA $\lam\!=\!0.1$ & 0.63\,\% & 0.27\,\% & 0.60\,\% \\
    \textbf{Exp9} & \textbf{DA $\lam\!=\!0.2$}
      & \textbf{0.58\,\%} & \textbf{0.27\,\%} & \textbf{0.58\,\%} \\
    \bottomrule
  \end{tabular}
\end{table}

\subsection{Dominant Error Patterns}

Table~\ref{tab:confusions} shows the top phoneme confusions.
Over 60\,\% of errors are vowel neutralizations inherent to
PT-BR~\cite{barbosa2004brazilian}.

\begin{table}[ht]
  \caption{Top phoneme confusions, 17.2M DA+dist.}
  \label{tab:confusions}
  \centering\small
  \begin{tabular}{lrl}
    \toprule
    Substitution & Count & Mechanism \\
    \midrule
    \textipa{E}$\to$e & 255 & Unstressed neutralization \\
    e$\to$\textipa{E} & 197 & Same (reverse) \\
    \textipa{O}$\to$o & 131 & Mid-back neutralization \\
    i$\to$e           & 121 & Vowel reduction \\
    o$\to$\textipa{O} & 95  & Mid-back neutralization \\
    \bottomrule
  \end{tabular}
\end{table}

Table~\ref{tab:voweldist} reveals the distributional asymmetry
driving the dominant confusion.
In pre-tonic position, /e/ outnumbers /\textipa{E}/ by 24.9:1; in
tonic position, the ratio inverts to 0.33:1.
The model learns the overwhelming pre-tonic bias and incorrectly
generalizes it to tonic syllables.

\begin{table}[ht]
  \caption{Corpus distribution of /e/ vs.\ /\textipa{E}/ by
           syllable position.}
  \label{tab:voweldist}
  \centering\small
  \begin{tabular}{lrrl}
    \toprule
    Position & /e/ & /\textipa{E}/ & Ratio \\
    \midrule
    Pre-tonic  & 38{,}147 & 1{,}535 & \textbf{24.9:1} \\
    Tonic      & 275      & 836     & 0.33:1 \\
    Post-tonic & 7{,}787  & 4{,}139 & 1.6:1 \\
    \midrule
    Global     & 46{,}209 & 6{,}510 & 7.1:1 \\
    \bottomrule
  \end{tabular}
\end{table}

DA~Loss cannot fix this: both /e/ and /\textipa{E}/ are
articulatorily close ($d\!=\!0.10$), so the domain penalty is
small regardless of direction.
This represents a principled limitation: domain-informed losses
help only when the error structure aligns with the distance metric.
Reducing these errors ($\sim$19\,\% of all substitutions) would
require explicit prosodic features or a dedicated tonic mid-vowel
classifier.

\subsection{Generalization Evidence}

Evaluation on 31~curated OOV words across 6~categories yields
55\,\% overall accuracy, with \textbf{100\,\%} on genuine novel
words absent from training (Table~\ref{tab:oov}).
Correct predictions include palatalization
(d+i$\to$d\textipa{Z}), coda reduction (l$\to$w), and nasal
mapping (om$\to$\~{o})---productive PT-BR phonological rules.
Defined failure modes: geminate consonants from Italian/English
loans (absent from training), English-phonology anglicisms
(genuinely OOV), and characters k/w/y (hard character-level OOV).

\begin{table}[ht]
  \caption{OOV generalization, 17.2M DA+dist model.}
  \label{tab:oov}
  \centering\small
  \begin{tabular}{lcc}
    \toprule
    Category & Correct & Phon.\ Score \\
    \midrule
    PT-BR neologisms        & 6/9\,(67\,\%)   & 97\,\% \\
    Geminate consonants     & 1/5\,(20\,\%)   & 81\,\% \\
    Anglicisms (in-vocab)   & 1/5\,(20\,\%)   & 71\,\% \\
    Char-OOV (k/w/y)        & 0/3\,(0\,\%)    & 68\,\% \\
    \textbf{Real PT-BR OOV} & \textbf{5/5\,(100\,\%)} & \textbf{100\,\%} \\
    Controls (in training)  & 4/4\,(100\,\%)  & 100\,\% \\
    \midrule
    \textbf{Total}          & \textbf{17/31\,(55\,\%)} & --- \\
    \bottomrule
  \end{tabular}
\end{table}

Table~\ref{tab:oov_examples} shows the 5~real OOV words with
exact IPA predictions---each exercising a different productive
PT-BR phonological rule.

\begin{table}[ht]
  \caption{Qualitative OOV predictions: 5/5 real novel PT-BR words.
           Rules: pal.\ = palatalization, red.\ = coda reduction.}
  \label{tab:oov_examples}
  \centering\small
  \begin{tabular}{lll}
    \toprule
    Word & Predicted IPA & Rule tested \\
    \midrule
    puxadinho & p u \textipa{S} a \textipa{"} d\textipa{Z} i \textipa{N} \textipa{U} & pal.\ d+i \\
    malcriado & m a w k \textipa{\*r} i \textipa{"} a d \textipa{U} & red.\ l$\to$w \\
    descartar & d\textipa{Z} i s k a x \textipa{"} t a x & coda x \\
    avermelhado & a v e x m e \textipa{"} \textipa{L} a d \textipa{U} & rr$\to$x \\
    computadorz\~{a}o & k \~{o} p u t a d o x \textipa{"} z \~{a} \~{\textipa{U}} & nasal mapping \\
    \bottomrule
  \end{tabular}
\end{table}

\subsection{Robustness Ablations}

Data reduction from 60\,\% to 50\,\% training degrades PER only
10\,\% relative (0.49\,$\to$\,0.54\,\%), suggesting the corpus is
above the saturation threshold.
Removing the hyphen character from the input vocabulary
(39$\to$38 tokens) has minimal impact (PER~0.58\,\%,
$+$0.04\,pp)---the hyphen is phonetically irrelevant for compound
words.

% =============================================================
\section{Discussion}
\label{sec:discussion}
% =============================================================

\textbf{When does domain-informed loss help?}
DA~Loss is most effective in the CE transition zone
(loss 0.3--1.5), where the model is actively resolving ambiguities
between candidate outputs.
Its bounded signal ($\leq \lam$) means it contributes $<$5\,\% of
gradient when the model is far from convergence.
This makes it a \emph{refinement signal}, not a replacement for CE.

\textbf{Interaction with representation.}
The non-additive interaction between DA~Loss and separators
(Table~\ref{tab:factorial}) suggests that output representation
determines whether the domain signal has ``room to operate.''
Without separators, DA differentiates phoneme confusions
effectively; with separators, the additional structural tokens
introduce a new error class (positional confusions) that the
distance metric addresses only partially.

\textbf{Generalizability.}
The framework (\ref{eq:diloss}) applies wherever a normalized
output distance exists: semantic similarity for NMT (WordNet,
embedding distance), acoustic distance for ASR, molecular
fingerprint distance for drug design.
The key requirement is that $d$ be \emph{domain-grounded}: random
or learned distances conflate the signal with the model's own
biases.

\textbf{Allophonic rule discovery.}
During OOV evaluation, the model produced /\textipa{G}/ in coda
before voiced consonants (e.g., \textit{borboleta} $\to$
b\,o\,\textipa{G}\,.\,b\,o\ldots).
Corpus audit confirmed the model was \textbf{correct}:
PT-BR voicing assimilation~\cite{barbosa2004brazilian} requires
/\textipa{G}/ before voiced consonants and /x/ in word-final
coda---perfect complementary distribution across all 95{,}937
dictionary entries with zero exceptions.
The model discovered this allophonic rule from distributional
evidence alone.

\textbf{Limitations.}
(1)~Evaluated on one task (G2P) and one language (PT-BR).
(2)~PanPhon's 24 binary features may underspecify some contrasts.
(3)~Perceptual validation (MOS/ABX) of ``closer errors are less
salient'' remains future work.
(4)~BiLSTM only; Transformer experiments are needed.

% =============================================================
\section{Conclusion}
% =============================================================

We presented a general framework for domain-informed loss functions
that augment cross-entropy with a distance penalty weighted by
prediction confidence.
Instantiated as Distance-Aware Loss for G2P with PanPhon
articulatory distance, the method reduces catastrophic errors by
19\,\% relative while achieving PER~0.48\,\% (Wilson CI
$\pm$0.03\,pp) on the largest reported PT-BR evaluation
(28{,}782~words).
A factorial analysis demonstrates that loss function design
interacts non-additively with model capacity and output
representation---motivating joint evaluation of these factors
in structured prediction tasks.
The framework is applicable to any domain providing a normalized
output distance metric.

% =============================================================
\section*{Use of Generative AI Tools}
% =============================================================

Generative AI tools (Claude, Anthropic) were used as editorial
assistants during manuscript preparation: grammar revision,
\LaTeX{} formatting, and cross-reference verification.
All scientific content---hypothesis formulation, experimental
design, data analysis, and interpretation of results---is the
sole work of the authors.
No AI-generated text appears in the scientific narrative without
human verification and revision.

% =============================================================
\bibliographystyle{IEEEtran}
\bibliography{mlsp_refs}
% =============================================================

\end{document}
